% theme: organisms_and_environment
\MajorQuestion{生物と環境}
\SetRowSpaceQanda

\Instruction{次の問いに答えなさい。}
\QQ{生物どうしの食べる・食べられるという関係が、鎖のようにつながったものを\NamedBlank{food_chain}という。}{食物連鎖}
\QQ{カビやキノコのなかまを\NamedBlank{fungi}という。}{菌類}
\QQ{生物と、それを取り巻く環境を一つのまとまりとしてとらえたものを\NamedBlank{ecosystem}という。}{生態系}
\QQ{乳酸菌や大腸菌などのなかまを\NamedBlank{bacteria}という。}{細菌類}
\QQ{光合成によって、無機物から有機物をつくり出す生物を\NamedBlank{producer}という。}{生産者}

\Instruction{\strut}
\QQ{人間の活動によって、ほかの地域から持ち込まれた生物を何というか。}{外来種（外来生物）}
\QQ{\RefBlank{fungi}の体をつくる、糸状の構造を何というか。}{菌糸}
\QQ{\RefBlank{producer}がつくった有機物を、直接または間接的に取り入れる生物を何というか。}{消費者}
\QQ{炭素や酸素などが形を変えながら\RefBlank{ecosystem}の中をめぐることを何というか。}{物質の循環}
\QQ{\RefBlank{food_chain}において、食べる・食べられる関係で分けた各段階を何というか。}{栄養段階}

\Instruction{\strut}
\QQ{生物の死がい等の有機物を、無機物に分解する過程に関わる生物を\NamedBlank{decomposer}という。}{分解者}
\QQ{地球全体の平均気温が長期的に上昇する現象を何というか。}{地球温暖化}
\QQ{\RefBlank{fungi}がつくり、なかまをふやすもとになるものを何というか。}{胞子}
\QQ{実際の自然界で、複数の\RefBlank{food_chain}が網の目のように複雑につながったものを何というか。}{食物網}
\QQ{一つの細胞だけで体ができている生物を何というか。}{単細胞生物}

\Instruction{\strut}
\QQ{自然環境を積極的に維持しようとすることを\NamedBlank{conservation}という。}{保全}
\QQ{ある物質の生物体内での濃度が外界より高くなり、\RefBlank{food_chain}を通してさらに高まる現象を\NamedBlank{bioaccumulation}という。}{生物濃縮}
\QQ{ある地域にもともと生息・生育していた生物を何というか。}{在来種（在来生物）}
\QQ{\RefBlank{bacteria}が一つの体を二つに分けてふえることを何というか。}{分裂}
\QQ{生物の種が地球上から完全にいなくなることを何というか。}{絶滅}
